\chapter{Deferred Proofs}

\vfill
\localtableofcontents
\vfill
\newpage

\section[Proof of Proposition 2.3]{Proof of \Cref{th:Kalman}}\label{sec:proof_Kalman}

\textbf{Prediction.} Let $\mathbb{V}[X]$ denote the variance of the random variable $X$:
\begin{equation}
	\mathbb{V}[X]\triangleq\mathbb{E}\left[(X-\mathbb{E}[X])(X-\mathbb{E}[X])^\top\right]
\end{equation}

Since $x_k=F_k\:x_{k-1}+b_k^{(e)}+w_k$ is an affine transformation of two independent Gaussian random variables $x_{k-1}|\mathbf{z}_{1:k-1}\sim\mathcal{N}\left(\hat{x}_{k-1|k-1},P_{k-1|k-1}\right)$ and $w_k\sim\mathcal{N}(0,Q_k)$, the predicted state $x_k|\mathbf{z}_{1:k-1}$ is also Gaussian, with:
\begin{equation}
	\begin{aligned}
	\hat{x}_{k|k-1}&\triangleq\mathbb{E}[x_k|\mathbf{z}_{1:k-1}]\\
	&=F_k\:\mathbb{E}[x_{k-1}|\mathbf{z}_{1:k-1}]+b_k^{(e)}+\mathbb{E}[w_k|\mathbf{z}_{1:k-1}]\\
	&=F_k\:\hat{x}_{k-1|k-1}+b_k^{(e)}
	\end{aligned}
\end{equation}
and:
\begin{equation}
	\begin{aligned}
	P_{k|k-1}&\triangleq\mathbb{V}[x_k|\mathbf{z}_{1:k-1}]\\
	&=F_k\:\mathbb{V}[x_{k-1}|\mathbf{z}_{1:k-1}]\:F_k^\top+\mathbb{V}[w_k|\mathbf{z}_{1:k-1}]\\
	&=F_k\:P_{k-1|k-1}\:F_k^\top+Q_k
	\end{aligned}
\end{equation}

The predicted covariance is thus SPD as the sum of an SPSD matrix and the SPD matrix $Q_k$.

\textbf{Correction.} The joint distribution of $(x_k,z_k)$ given $\mathbf{z}_{1:k-1}$ is also Gaussian, since $z_k=H_k\:x_k+b_k^{(o)}+n_k$ is an affine transformation of the independent Gaussian variables $x_k|\mathbf{z}_{1:k-1}\sim\mathcal{N}\left(\hat{x}_{k|k-1},P_{k|k-1}\right)$ and $n_k\sim\mathcal{N}(0,R_k)$. It comes that:
\begin{equation}
	\left.\begin{bmatrix}
	x_k\\z_k
	\end{bmatrix}\right|
	\mathbf{z}_{1:k-1}\sim\mathcal{N}
	\left(\begin{bmatrix}
	\hat{x}_{k|k-1}\\H_k\hat{x}_{k|k-1}+b_k^{(o)}
	\end{bmatrix},
	\begin{bmatrix}
	P_{k|k-1}&P_{k|k-1}H_k^\top\\
	H_k P_{k|k-1}&S_k
	\end{bmatrix}\right)
\end{equation}
where:
\begin{equation}
	S_k\triangleq H_k\:P_{k|k-1}\:H_k^\top+R_k
\end{equation}
is called the innovation covariance, and is the covariance associated with the innovation vector:
\begin{equation}
	y_k\triangleq z_k-\left(H_k\:\hat{x}_{k|k-1}+b_k^{(o)}\right)
\end{equation}

Since $R_k\succ0$, the independent pair $(x_k,n_k)$ has the covariance:
\begin{equation}
	\begin{bmatrix}
	P_{k|k-1}&0\\0&R_k
	\end{bmatrix}\succ0
\end{equation}

Its image under the invertible linear map $(x,n)\mapsto(x,H_k\:x+n)$ therefore has an SPD covariance as well. The Gaussian conditioning formula (see \Cref{th:Gaussian_conditioning}) can thus be applied to give the posterior distribution $x_k|\mathbf{z}_{1:k}$ is Gaussian, with:
\begin{align}
	\hat{x}_{k|k}&=\hat{x}_{k|k-1}+K_k\:y_k\\
	P_{k|k}&=(I-K_k\:H_k)\:P_{k|k-1}
\end{align}
where:
\begin{equation}
	K_k\triangleq P_{k|k-1}\:H_k^\top\:S_k^{-1}
\end{equation}
is the Kalman gain.

Finally, \Cref{th:Gaussian_conditioning} proves that the Schur complement of the SPD joint covariance:
\begin{equation}
	P_{k|k-1}-\left(P_{k|k-1}\:H_k^\top\right)S_k^{-1}\left(H_k\:P_{k|k-1}\right)=P_{k|k-1}-K_k\:H_k\:P_{k|k-1}=P_{k|k}
\end{equation}
is itself SPD, so $P_{k|k}\succ0$.

\section[Proof of Proposition 2.10]{Proof of \Cref{th:RBPF_decorrelation}}\label{sec:proof_RBPF}

As a linear combination of the centred noises $w_k^l$ and $w_k^n$, $\overline{w}_k^l$ is centred. Its covariance develops as:
\begin{equation}
	\begin{aligned}
	\overline{Q}_k^l&=\mathbb{E}\left[\overline{w}_k^l\:{\overline{w}_k^l}^\top\right]\\
	&=\mathbb{E}\left[(w_k^l-Q_k^{nl}\:(Q_k^n)^{-1}\:w_k^n)\:(w_k^l-Q_k^{nl}\:(Q_k^n)^{-1}\:w_k^n)^\top\right]\\
	&=\mathbb{E}\left[w_k^l\:(w_k^l)^\top\right]-\mathbb{E}\left[w_k^l\:(Q_k^{nl}\:(Q_k^n)^{-1}\:w_k^n)^\top\right]-\mathbb{E}\left[(Q_k^{nl}\:(Q_k^n)^{-1}\:w_k^n)\:(w_k^l)^\top\right]\\
	&\qquad+\mathbb{E}\left[(Q_k^{nl}\:(Q_k^n)^{-1}\:w_k^n)\:(Q_k^{nl}\:(Q_k^n)^{-1}\:w_k^n)^\top\right]\\
	&=\mathbb{E}\left[w_k^l\:(w_k^l)^\top\right]-\mathbb{E}\left[w_k^l\:(w_k^n)^\top\right](Q_k^n)^{-\top}\:(Q_k^{nl})^\top-Q_k^{nl}\:(Q_k^n)^{-1}\:\mathbb{E}\left[w_k^n\:(w_k^l)^\top\right]\\
	&\qquad+Q_k^{nl}\:(Q_k^n)^{-1}\:\mathbb{E}\left[w_k^n\:(w_k^n)^\top\right]\:(Q_k^n)^{-\top}\:(Q_k^{nl})^\top\\
	&=Q_k^l-Q_k^{nl}\:(Q_k^n)^{-\top}\:(Q_k^{nl})^\top-Q_k^{nl}\:(Q_k^n)^{-1}\:(Q_k^{nl})^\top+Q_k^{nl}\:(Q_k^n)^{-1}\:Q_k^n\:(Q_k^n)^{-\top}\:(Q_k^{nl})^\top\\
	&=Q_k^l-Q_k^{nl}\:(Q_k^n)^{-1}\:(Q_k^{nl})^\top-Q_k^{nl}\:(Q_k^n)^{-1}\:(Q_k^{nl})^\top+Q_k^{nl}\:(Q_k^n)^{-1}\:(Q_k^{nl})^\top\\
	&=Q_k^l-Q_k^{nl}\:(Q_k^n)^{-1}\:(Q_k^{nl})^\top
	\end{aligned}
\end{equation}

Finally, computing:
\begin{equation}
	\begin{aligned}
	\mathbb{E}\left[\overline{w}_k^l\:(w_k^n)^\top\right]&=\mathbb{E}\left[(w_k^l-Q_k^{nl}\:(Q_k^n)^{-1}\:w_k^n)\:(w_k^n)^\top\right]\\
	&=\mathbb{E}\left[w_k^l\:(w_k^n)^\top\right]-Q_k^{nl}\:(Q_k^n)^{-1}\:\mathbb{E}\left[w_k^n\:(w_k^n)^\top\right]\\
	&=Q_k^{nl}-Q_k^{nl}\:(Q_k^n)^{-1}\:Q_k^n\\
	&=0
	\end{aligned}
\end{equation}
shows that $\overline{w}_k^l$ and $w_k^n$ are uncorrelated.

\section[Proof of Proposition 3.2]{Proof of \Cref{th:kernel_gram_equivalence}}\label{sec:proof_definite_kernels}

\textbf{Forward implication.} Let $n\in\mathbb{N}^*$ and $\mathbf{x}\in\mathcal{U}^n$. It comes from the symmetry of $k$ that:
\begin{equation}
	\begin{aligned}
	k(\mathbf{x},\mathbf{x})&=\begin{bmatrix}
	k(x_1,x_1)&\cdots&k(x_1,x_n)\\\vdots&\ddots&\vdots\\k(x_n,x_1)&\cdots&k(x_n,x_n)
	\end{bmatrix}\\
	&=\begin{bmatrix}
	k(x_1,x_1)^\top&\cdots&k(x_n,x_1)^\top\\\vdots&\ddots&\vdots\\k(x_1,x_n)^\top&\cdots&k(x_n,x_n)^\top
	\end{bmatrix}\\
	&=\begin{bmatrix}
	k(x_1,x_1)&\cdots&k(x_1,x_n)\\\vdots&\ddots&\vdots\\k(x_n,x_1)&\cdots&k(x_n,x_n)
	\end{bmatrix}^\top\\
	&=k(\mathbf{x},\mathbf{x})^\top
	\end{aligned}
\end{equation}

Let $a\in\mathbb{R}^{nd_z}$. We decompose $a=[a_1^\top,\dots,a_n^\top]^\top$ with $a_i\in\mathbb{R}^{d_z},\:i\in[\![1,n]\!]$ to obtain:
\begin{equation}
	\begin{aligned}
	a^\top k(\mathbf{x},\mathbf{x})\:a&=\begin{bmatrix}a_1^\top&\cdots&a_n^\top\end{bmatrix}\begin{bmatrix}k(x_1,x_1)&\cdots&k(x_1,x_n)\\\vdots&\ddots&\vdots\\k(x_n,x_1)&\cdots&k(x_n,x_n)\end{bmatrix}\begin{bmatrix}a_1\\\vdots\\a_n\end{bmatrix}\\
	&=\sum_{i=1}^n\sum_{j=1}^na_i^\top k(x_i,x_j)\:a_j\\
	&\geq0
	\end{aligned}
\end{equation}

\textbf{Reverse implication.} Let $x,x'\in\mathcal{U}$. By symmetry of $k((x,x'),(x,x'))$, it comes that:
\begin{equation}
	\begin{aligned}
	\begin{bmatrix}
	k(x,x)&k(x,x')\\k(x',x)&k(x',x')
	\end{bmatrix}&=\begin{bmatrix}
	k(x,x)&k(x,x')\\k(x',x)&k(x',x')
	\end{bmatrix}^\top\\
	&=\begin{bmatrix}
	k(x,x)^\top&k(x',x)^\top\\k(x,x')^\top&k(x',x')^\top
	\end{bmatrix}
	\end{aligned}
\end{equation}

In particular, the upper-right block gives:
\begin{equation}
	k(x,x')=k(x',x)^\top
\end{equation}

Let $n\in\mathbb{N}^*$, $\mathbf{x}\in\mathcal{U}^n$, and $a_1,\dots,a_n\in\mathbb{R}^{d_z}$, and set $a\triangleq[a_1^\top,\dots,a_n^\top]^\top\in\mathbb{R}^{nd_z}$. Then:
\begin{equation}
	\begin{aligned}
	\sum_{i=1}^n\sum_{j=1}^na_i^\top k(x_i,x_j)\:a_j&=\begin{bmatrix}a_1^\top&\cdots&a_n^\top\end{bmatrix}\begin{bmatrix}k(x_1,x_1)&\cdots&k(x_1,x_n)\\\vdots&\ddots&\vdots\\k(x_n,x_1)&\cdots&k(x_n,x_n)\end{bmatrix}\begin{bmatrix}a_1\\\vdots\\a_n\end{bmatrix}\\
	&=a^\top k(\mathbf{x},\mathbf{x})\:a\\
	&\geq0
	\end{aligned}
\end{equation}

\textbf{Strict version.} The same computations hold with strict inequalities, provided the inputs are pairwise distinct and the coefficients not all zero.

\section[Proof of Proposition 3.3]{Proof of \Cref{th:Matérn_properties}}\label{sec:proof_Matérn}

Assume that $\nu\in\mathbb{R}_+^*$ and let $k_\nu$ denote such a function. We also let:
\begin{equation}
	g:\left\{\begin{aligned}
	\mathbb{R}_+&\to\mathbb{R}_+\\
	r&\mapsto\sigma_f^2\:\frac{2^{1-\nu}}{\Gamma(\nu)}r^\nu K_\nu(r)
	\end{aligned}\right.
	\quad\text{and}\quad
	\tilde{r}:\left\{\begin{aligned}
	\mathbb{R}^{d_x}\times\mathbb{R}^{d_x}&\to\mathbb{R}_+\\
	x,x'&\mapsto\frac{\sqrt{2\nu}}{l}\|x-x'\|
	\end{aligned}\right.
\end{equation}

$\tilde{r}$ and $K_\nu$ are continuous on $\mathbb{R}^{d_x}\times\mathbb{R}^{d_x}$ and $\mathbb{R}_+^*$ respectively, and the latter admits the following asymptotic expansion as $r\to0^+$:
\begin{equation}
	K_\nu(r)=2^{\nu-1}\Gamma(\nu)\:r^{-\nu}+o\left(r^{-\nu}\right)
\end{equation}

Therefore, $f$ can be continuously extended to $0$ by setting $f(0)=\sigma_f^2$, and $k_\nu$ is thus continuous on the entirety of $\mathbb{R}^{d_x}\times\mathbb{R}^{d_x}$.

Let $x,x'\in\mathbb{R}^{d_x}$. By symmetry of the norm, $\tilde{r}(x,x')=\tilde{r}(x',x)$, and thus:
\begin{equation}
	k_\nu=g\circ\tilde{r}
\end{equation}
is symmetric, with $\circ$ being the composition operator.

Furthermore, Bochner's theorem \cite{bib:Akhiezer} states that a continuous integrable function with a strictly positive Fourier transform is SPD. It can be shown \cite{bib:Matérn-correlation} that the Fourier transform of $k_\nu$ on $\mathbb{R}^{d_x}$ is:
\begin{equation}
	S_\nu:\left\{\begin{aligned}
	\mathbb{R}^{d_x}&\to\mathbb{R}\\
	\omega&\mapsto C\:\left(\frac{2\nu}{l^2}+\|\omega\|^2\right)^{-(\nu+d_x/2)}
	\end{aligned}\right.
\end{equation}
where $C>0$ is a normalising constant which only depends on $d_x$ and the parameters of $k_\nu$. Since $\nu>0$, $l>0$, it comes that:
\begin{equation}
	\forall\omega\in\mathbb{R}^{d_x},\:S_\nu(\omega)>0
\end{equation}
and thus $k_\nu$ is a SPD kernel. By definition, it is also a radial kernel.

The case $\nu=+\infty$ is more direct. The squared exponential kernel is continuous on the whole of $\mathbb{R}_+$ without any extension, and symmetric by the same composition argument. Its Fourier transform on $\mathbb{R}^{d_x}$ is:
\begin{equation}
	S_\nu:\left\{\begin{aligned}
	\mathbb{R}^{d_x}&\to\mathbb{R}\\
	\omega&\mapsto C\:\exp\left(-\frac{l^2\|\omega\|^2}{2}\right)
	\end{aligned}\right.
\end{equation}
which is again strictly positive, so Bochner's theorem applies as above.

\section[Proof of Proposition 3.8]{Proof of \Cref{th:ordinary_kriging}}\label{sec:proof_OK}

Let:
\begin{equation}
	\mathcal{L}(W,\xi)\triangleq\operatorname{tr}\left(\mathbb{V}\left[W^\top\tilde{\mathbf{z}}-z\right]\right)+2\:\operatorname{tr}\left(\xi^\top\left(C^\top W-I_{d_z}\right)\right)
\end{equation}
be the Lagrangian associated with the optimisation problem \eqref{eq:error_covariance} under the constraint \eqref{eq:ordinary_kriging_constraint}, and $\xi\in\mathcal{M}_{d_z,d_z}(\mathbb{R})$ the corresponding Lagrange multiplier.

This problem is solved by finding the minimum of $\mathcal{L}$, \textit{i.e.} finding a zero of its derivative. Writing $G$ for $G(\tilde{\mathbf{x}})$, the same properties used in the proof of \Cref{th:simple_kriging} give:
\begin{equation}
	\begin{aligned}
	&\frac{\partial\mathcal{L}}{\partial W}(W,\xi)=0\\
	\iff\quad&2G\:W-2k(\tilde{\mathbf{x}},x)+2\frac{\partial}{\partial W}\left(\operatorname{tr}\left(\xi^\top\left(C^\top W-I_{d_z}\right)\right)\right)=0\\
	\iff\quad&2G\:W-2k(\tilde{\mathbf{x}},x)+2C\xi=0\\
	\iff\quad&G\:W+C\xi=k(\tilde{\mathbf{x}},x)
	\end{aligned}\label{eq:Lagrangian_minimum}
\end{equation}

Together with the constraint $C^\top W=I_{d_z}$, this yields the block system:
\begin{equation}
	\begin{bmatrix}
	G&C\\C^\top&0
	\end{bmatrix}\begin{bmatrix}
	W\\\xi
	\end{bmatrix}=\begin{bmatrix}
	k(\tilde{\mathbf{x}},x)\\I_{d_z}
	\end{bmatrix}
\end{equation}

The left matrix block can be inverted using the upper left Schur inversion formula from \Cref{th:Schur}. Its Schur complement,
\begin{equation}
	S\triangleq0-C^\top G^{-1}C=-C^\top G^{-1}C,
\end{equation}
is invertible since $G\succ0$ and $C$ has full column rank. The convention of \Cref{th:Schur} gives the four blocks of the inverse:
\begin{equation}
	\left\{\begin{aligned}
	\bar{A}&=G^{-1}+G^{-1}C\:S^{-1}C^\top G^{-1}\\
	\bar{B}&=-G^{-1}C\:S^{-1}\\
	\bar{C}&=-S^{-1}C^\top G^{-1}\\
	\bar{D}&=S^{-1}
	\end{aligned}\right.
\end{equation}

Let $M\triangleq-S^{-1}=\left(C^\top G^{-1}C\right)^{-1}$ and $U\triangleq I_{d_z}-C^\top G^{-1}k(\tilde{\mathbf{x}},x)$. The weight matrix can be derived from the upper blocks of the inverse:
\begin{equation}
	\begin{aligned}
	W&=\bar{A}\:k(\tilde{\mathbf{x}},x)+\bar{B}\:I_{d_z}\\
	&=\left(G^{-1}+G^{-1}C\:S^{-1}C^\top G^{-1}\right)k(\tilde{\mathbf{x}},x)-G^{-1}C\:S^{-1}\\
	&=G^{-1}k(\tilde{\mathbf{x}},x)+G^{-1}C\:S^{-1}\left(C^\top G^{-1}k(\tilde{\mathbf{x}},x)-I_{d_z}\right)\\
	&=G^{-1}k(\tilde{\mathbf{x}},x)+G^{-1}C\:M\left(I_{d_z}-C^\top G^{-1}k(\tilde{\mathbf{x}},x)\right)\\
	&=G^{-1}k(\tilde{\mathbf{x}},x)+G^{-1}C\:M\:U
	\end{aligned}
\end{equation}
and $\xi$ from the lower blocks:
\begin{equation}
	\begin{aligned}
	\xi&=\bar{C}\:k(\tilde{\mathbf{x}},x)+\bar{D}\:I_{d_z}\\
	&=-S^{-1}C^\top G^{-1}k(\tilde{\mathbf{x}},x)+S^{-1}\\
	&=S^{-1}\left(I_{d_z}-C^\top G^{-1}k(\tilde{\mathbf{x}},x)\right)\\
	&=-M\left(I_{d_z}-C^\top G^{-1}k(\tilde{\mathbf{x}},x)\right)\\
	&=-M\:U
	\end{aligned}
\end{equation}

Since $M$ is symmetric, the OK estimator is:
\begin{equation}
	\begin{aligned}
	\hat{z}_{\mathrm{OK}}&=W^\top\tilde{\mathbf{z}}\\
	&=k(x,\tilde{\mathbf{x}})\:G^{-1}\tilde{\mathbf{z}}+U^\top M\:C^\top G^{-1}\tilde{\mathbf{z}}\\
	&=k(x,\tilde{\mathbf{x}})\:G^{-1}\tilde{\mathbf{z}}+\left(I_{d_z}-k(x,\tilde{\mathbf{x}})\:G^{-1}C\right)M\:C^\top G^{-1}\tilde{\mathbf{z}}\\
	&=M\:C^\top G^{-1}\tilde{\mathbf{z}}+k(x,\tilde{\mathbf{x}})\:G^{-1}\left(\tilde{\mathbf{z}}-C\:M\:C^\top G^{-1}\tilde{\mathbf{z}}\right)
	\end{aligned}
\end{equation}
which is exactly \eqref{eq:ordinary_kriging_equations}, the estimated mean being $\hat{m}_{\mathrm{OK}}(\tilde{\mathbf{x}},\tilde{\mathbf{z}})=M\:C^\top G^{-1}\tilde{\mathbf{z}}$.

Since $R_{\mathrm{OK}}$ is a centred moment of the estimation error, it is insensitive to the unknown mean and \eqref{eq:generic_kriging_variance} can be applied. Using \eqref{eq:Lagrangian_minimum} and the constraint $W^\top C=I_{d_z}$ gives the identity:
\begin{equation}
	\begin{aligned}
	W^\top G\:W&=W^\top\left(k(\tilde{\mathbf{x}},x)-C\:\xi\right)\\
	&=W^\top k(\tilde{\mathbf{x}},x)-W^\top C\:\xi\\
	&=W^\top k(\tilde{\mathbf{x}},x)-\xi
	\end{aligned}\label{eq:OK_identity}
\end{equation}

Combining these results yields:
\begin{equation}
	\begin{aligned}
	R_{\mathrm{OK}}&=W^\top G\:W-W^\top k(\tilde{\mathbf{x}},x)-k(x,\tilde{\mathbf{x}})\:W+k(x,x)+R\\
	&=W^\top k(\tilde{\mathbf{x}},x)-\xi-W^\top k(\tilde{\mathbf{x}},x)-k(x,\tilde{\mathbf{x}})\:W+k(x,x)+R\\
	&=-\xi-k(x,\tilde{\mathbf{x}})\:W+k(x,x)+R\\
	&=M\:U-k(x,\tilde{\mathbf{x}})\:G^{-1}k(\tilde{\mathbf{x}},x)-k(x,\tilde{\mathbf{x}})\:G^{-1}C\:M\:U+k(x,x)+R\\
	&=\left(I_{d_z}-k(x,\tilde{\mathbf{x}})\:G^{-1}C\right)M\:U+k(x,x)+R-k(x,\tilde{\mathbf{x}})\:G^{-1}k(\tilde{\mathbf{x}},x)\\
	&=U^\top M\:U+k(x,x)+R-k(x,\tilde{\mathbf{x}})\:G^{-1}k(\tilde{\mathbf{x}},x)
	\end{aligned}
\end{equation}

Finally, recognising $R_{\hat{m}_{\mathrm{OK}}}(x,\tilde{\mathbf{x}})=U^\top M\:U\succeq0$, since $M\succ0$, gives the desired expression.

\section[Proof of Proposition 3.9]{Proof of \Cref{th:universal_kriging}}\label{sec:proof_UK}

The unbiased assumption now reads:
\begin{equation}
	g(\tilde{\mathbf{x}})^\top W-g(x)^\top=0
\end{equation}
and the Lagrange multiplier method of \Cref{th:ordinary_kriging} can be transposed to universal kriging under the substitutions:
\begin{equation}
	\left\{\begin{aligned}
	C&\longrightarrow g(\tilde{\mathbf{x}})\\
	I_{d_z}&\longrightarrow g(x)^\top\\
	\xi&\in\mathcal{M}_{r_{\mathrm{UK}}d_z,d_z}(\mathbb{R})
	\end{aligned}\right.
\end{equation}

The full column rank of $g(\tilde{\mathbf{x}})$ ensures that $S$ is invertible, and that $M$ is SPD. Writing $G$ for $G(\tilde{\mathbf{x}})$, the shorthands of that proof thus become:
\begin{equation}
	\left\{\begin{aligned}
	M&\triangleq-S^{-1}=\left(g(\tilde{\mathbf{x}})^\top G^{-1}g(\tilde{\mathbf{x}})\right)^{-1}\\
	U&\triangleq g(x)^\top-g(\tilde{\mathbf{x}})^\top G^{-1}k(\tilde{\mathbf{x}},x)
	\end{aligned}\right.
\end{equation}
so that the weight matrix and Lagrange multiplier are:
\begin{equation}
	\left\{\begin{aligned}
	W&=G^{-1}k(\tilde{\mathbf{x}},x)+G^{-1}g(\tilde{\mathbf{x}})\:M\:U\\
	\xi&=-M\:U
	\end{aligned}\right.
\end{equation}

Finally, the bilinear identity \eqref{eq:OK_identity} becomes:
\begin{equation}
	W^\top G\:W=W^\top k(\tilde{\mathbf{x}},x)-g(x)\:\xi
\end{equation}

The UK estimator and covariance matrices are derived exactly like for \Cref{th:ordinary_kriging}, using the corrected results:
\begin{equation}
	\begin{aligned}
	\hat{z}_{\mathrm{UK}}&=k(x,\tilde{\mathbf{x}})\:G^{-1}\tilde{\mathbf{z}}+U^\top M\:g(\tilde{\mathbf{x}})^\top G^{-1}\tilde{\mathbf{z}}\\
	&=g(x)\:M\:g(\tilde{\mathbf{x}})^\top G^{-1}\tilde{\mathbf{z}}+k(x,\tilde{\mathbf{x}})\:G^{-1}\left(\tilde{\mathbf{z}}-g(\tilde{\mathbf{x}})\:M\:g(\tilde{\mathbf{x}})^\top G^{-1}\tilde{\mathbf{z}}\right)
	\end{aligned}\label{eq:UK_z_proof}
\end{equation}
and:
\begin{equation}
	R_{\mathrm{UK}}=U^\top M\:U+k(x,x)+R-k(x,\tilde{\mathbf{x}})\:G^{-1}k(\tilde{\mathbf{x}},x)\label{eq:UK_R_proof}
\end{equation}

\section[Proof of Proposition 3.10]{Proof of \Cref{th:kriging_vague_prior}}\label{sec:proof_vague_prior}

Using \eqref{eq:UK_mean}, it comes that:
\begin{equation}
	f\mid\beta\sim\mathcal{GP}(g\:\beta,k)
\end{equation}
and thus:
\begin{equation}
	f-g\:\beta\mid\beta\sim\mathcal{GP}(0,k)
\end{equation}

Since this law is free of $\beta$, $f-g\:\beta$ is independent of $\beta$, which gives:
\begin{equation}
	f-g\:\beta\sim\mathcal{GP}(0,k)
\end{equation}

Let $x,x'\in\mathcal{U}$. The mean of $f$ can thus be computed by decomposing:
\begin{equation}
	\mathbb{E}[f(x)]=g(x)\:\mathbb{E}[\beta]+\mathbb{E}[f(x)-g(x)\:\beta]=0
\end{equation}

By independence of $f-g\:\beta$ and $\beta$, it follows that:
\begin{equation}
	\begin{aligned}
	\mathbb{E}\left[f(x)\:f(x')^\top\right]&=\mathbb{E}\left[(f(x)-g(x)\:\beta)(f(x')-g(x')\:\beta)^\top\right]+g(x)^\top\:\mathbb{E}\left[\beta\:\beta^\top\right]g(x')+\\
	&\qquad\mathbb{E}\left[(f(x)-g(x)\:\beta)\:\beta^\top\right]g(x')^\top+g(x')\:\mathbb{E}\left[\beta\:(f(x)-g(x)\:\beta)^\top\right]\\
	&=k(x,x')+g(x)\:B\:g(x')^\top+0+0\\
	&=k_B(x,x')
	\end{aligned}
\end{equation}

As a linear combiantion of the jointly Gaussian distributions $f-g\:\beta$ and $\beta$, we finally have:
\begin{equation}
	f\sim\mathcal{GP}(0,k_B)
\end{equation}

In this form, the mean function of $f$ is known, and \Cref{th:simple_kriging_conditional} states that simple kriging using $k_B$ gives the conditional distribution of $z\mid\tilde{\mathbf{z}}$:
\begin{equation}
	z\mid\tilde{\mathbf{z}}\sim\mathcal{N}\left(\hat{z}_{\mathrm{SK},B}(x,\tilde{\mathbf{x}},\tilde{\mathbf{z}},0),R_{\mathrm{SK},B}(x,\tilde{\mathbf{x}})\right)
\end{equation}

The SK estimator and error covariance matrix can be expressed as follows:
\begin{equation}
	\left\{\begin{aligned}
	\hat{z}_{\mathrm{SK},B}(x,\tilde{\mathbf{x}},\tilde{\mathbf{z}},0)&=k_B(x,\tilde{\mathbf{x}})\:G_B(\tilde{\mathbf{x}})^{-1}\:\tilde{\mathbf{z}}\\
	R_{\mathrm{SK},B}(x,\tilde{\mathbf{x}})&=k_B(x,x)+R-k_B(x,\tilde{\mathbf{x}})\:G_B(\tilde{\mathbf{x}})^{-1}k_B\left(\tilde{\mathbf{x}},x\right)
	\end{aligned}\right.
\end{equation}
where:
\begin{equation}
	G_B\triangleq k_B(\tilde{\mathbf{x}},\tilde{\mathbf{x}})+I_{\tilde{N}}\otimes\tilde{R}=G+g(\tilde{\mathbf{x}})\:B\:g(\tilde{\mathbf{x}})^\top
\end{equation}

For the sake of clarity, we let $G$ and $g$ respectively denote $G(\tilde{\mathbf{x}})$ and $g(\tilde{\mathbf{x}})$, and we reuse:
\begin{equation}
	\left\{\begin{aligned}
	M&\triangleq\left(g^\top G^{-1}g\right)^{-1}\\
	U&\triangleq g(x)^\top-g^\top G^{-1}k(\tilde{\mathbf{x}},x)
	\end{aligned}\right.
\end{equation}
from \cref{sec:proof_UK}.

Using \Cref{th:Woodbury}, it comes that:
\begin{equation}
	G_B^{-1}=G^{-1}-G^{-1}g\:M_B\:g^\top G^{-1}
\end{equation}
with:
\begin{equation}
	M_B\triangleq\left(B^{-1}+g^\top G^{-1}g\right)^{-1}=\left(B^{-1}+M^{-1}\right)^{-1}
\end{equation}

By definition of $M_B$, we have:
\begin{equation}
	\begin{aligned}
	&\left(B^{-1}+M^{-1}\right)M_B=I\\
	\iff\quad&M^{-1}M_B=I-B^{-1}M_B\\
	\iff\quad&B\:M^{-1}M_B=B-M_B\\
	\end{aligned}
\end{equation}

We also precompute:
\begin{equation}
	\begin{aligned}
	g(x)\:B\:g^\top\:G_B^{-1}&=g(x)\:B\:g^\top\left(G^{-1}-G^{-1}g\:M_B\:g^\top G^{-1}\right)\\
	&=g(x)\left(B\:g^\top G^{-1}-B\:g^\top G^{-1}g\:M_B\:g^\top G^{-1}\right)\\
	&=g(x)\left(B\:g^\top G^{-1}-B\:M^{-1}\:M_B\:g^\top G^{-1}\right)\\
	&=g(x)\left(B\:g^\top G^{-1}-(B-M_B)\:g^\top G^{-1}\right)\\
	&=g(x)\:M_B\:g^\top G^{-1}
	\end{aligned}
\end{equation}

It comes that:
\begin{equation}
	\begin{aligned}
	k_B(x,\tilde{\mathbf{x}})\:G_B^{-1}&=\left(k(x,\tilde{\mathbf{x}})+g(x)\:B\:g^\top\right)G_B^{-1}\\
	&=k(x,\tilde{\mathbf{x}})\:G_B^{-1}+g(x)\:B\:g^\top G_B^{-1}\\
	&=k(x,\tilde{\mathbf{x}})\left(G^{-1}-G^{-1}g\:M_B\:g^\top G^{-1}\right)+g(x)\:M_B\:g^\top G^{-1}\\
	&=k(x,\tilde{\mathbf{x}})\:G^{-1}+\left(-k(x,\tilde{\mathbf{x}})\:G^{-1}g+g(x)\right)M_B\:g^\top G^{-1}\\
	&=k(x,\tilde{\mathbf{x}})\:G^{-1}+U^\top M_B\:g^\top G^{-1}
	\end{aligned}
\end{equation}
and thus:
\begin{equation}
	\hat{z}_{\mathrm{SK},B}=k_B(x,\tilde{\mathbf{x}})\:G_B^{-1}\tilde{\mathbf{z}}=k(x,\tilde{\mathbf{x}})\:G^{-1}\tilde{\mathbf{z}}+U^\top M_B\:g^\top G^{-1}\tilde{\mathbf{z}}
\end{equation}

For the covariance, we compute:
\begin{equation}
	\begin{aligned}
	k_B(x,\tilde{\mathbf{x}})\:G_B^{-1}k_B(\tilde{\mathbf{x}},x)&=\left(k(x,\tilde{\mathbf{x}})\:G^{-1}+U^\top M_B\:g^\top G^{-1}\right)\left(k(\tilde{\mathbf{x}},x)+g\:B\:g(x)^\top\right)\\
	&=k(x,\tilde{\mathbf{x}})\:G^{-1}k(\tilde{\mathbf{x}},x)+U^\top M_B\:g^\top G^{-1}k(\tilde{\mathbf{x}},x)+\\
	&\qquad k(x,\tilde{\mathbf{x}})\:G^{-1}g\:B\:g(x)^\top+U^\top M_B\:g^\top G^{-1}g\:B\:g(x)^\top\\
	&=k(x,\tilde{\mathbf{x}})\:G^{-1}k(\tilde{\mathbf{x}},x)+U^\top M_B\:g^\top G^{-1}k(\tilde{\mathbf{x}},x)+\\
	&\qquad k(x,\tilde{\mathbf{x}})\:G^{-1}g\:B\:g(x)^\top+U^\top M_B\:M^{-1}\:B\:g(x)^\top\\
	&=k(x,\tilde{\mathbf{x}})\:G^{-1}k(\tilde{\mathbf{x}},x)+U^\top M_B\:g^\top G^{-1}k(\tilde{\mathbf{x}},x)+\\
	&\qquad k(x,\tilde{\mathbf{x}})\:G^{-1}g\:B\:g(x)^\top+U^\top(B-M_B)\:g(x)^\top
	\end{aligned}
\end{equation}
and:
\begin{equation}
	\begin{aligned}
	g(x)\:B\:g(x)^\top-k(x,\tilde{\mathbf{x}})\:G^{-1}g\:B\:g(x)^\top&=U^\top g\:B\:g(x)^\top
	\end{aligned}
\end{equation}
to obtain:
\begin{equation}
	\begin{aligned}
	R_{\mathrm{SK},B}&=k_B(x,x)+R-k_B(x,\tilde{\mathbf{x}})\:G_B(\tilde{\mathbf{x}})^{-1}k_B\left(\tilde{\mathbf{x}},x\right)\\
	&=k(x,x)+g(x)\:B\:g(x)^\top+R-k(x,\tilde{\mathbf{x}})\:G^{-1}k(\tilde{\mathbf{x}},x)-U^\top M_B\:g^\top G^{-1}k(\tilde{\mathbf{x}},x)-\\
	&\qquad k(x,\tilde{\mathbf{x}})\:G^{-1}g\:B\:g(x)^\top-U^\top(B-M_B)\:g(x)^\top\\
	&=k(x,x)+R-k(x,\tilde{\mathbf{x}})\:G^{-1}k(\tilde{\mathbf{x}},x)-U^\top M_B\:g^\top G^{-1}k(\tilde{\mathbf{x}},x)+U^\top M_B\:g(x)^\top\\
	&=k(x,x)+R-k(x,\tilde{\mathbf{x}})\:G^{-1}k(\tilde{\mathbf{x}},x)+U^\top M_B\:U
	\end{aligned}
\end{equation}

Since $M_B\to M$ as $B^{-1}\to0$, we finally have:
\begin{equation}
	\left\{\begin{aligned}
	\hat{z}_{\mathrm{SK},B}(x,\tilde{\mathbf{x}},\tilde{\mathbf{z}},0)&\underset{B^{-1}\to0}{\longrightarrow}k(x,\tilde{\mathbf{x}})\:G^{-1}\tilde{\mathbf{z}}+U^\top M\:g^\top G^{-1}\tilde{\mathbf{z}}\\
	R_{\mathrm{SK},B}(x,\tilde{\mathbf{x}})&\underset{B^{-1}\to0}{\longrightarrow}k(x,x)+R-k(x,\tilde{\mathbf{x}})\:G^{-1}k(\tilde{\mathbf{x}},x)+U^\top M\:U
	\end{aligned}\right.
\end{equation}
the limit of which are exactly \eqref{eq:UK_z_proof} and \eqref{eq:UK_R_proof}.

\begin{comment}
\textbf{TODO: ça a l'air de mériter sa propre proposition mais jsp si ç'est utile}The coefficients follow the same course. The Gaussian linear model gives $\beta$ the posterior $\mathcal{N}\left(M_B\:g^\top G^{-1}\tilde{\mathbf{z}},M_B\right)$, whose limit is $\mathcal{N}\left(\hat{m}_{\mathrm{UK}}(\tilde{\mathbf{x}},\tilde{\mathbf{z}},g),M\right)$. The estimator of \eqref{eq:kriged_universal_mean} is therefore the posterior mean of the trend, and $M$ its posterior covariance.
\end{comment}
